\documentclass[12pt,a4paper]{report}

% ── Packages ────────────────────────────────────────────────
\usepackage[utf8]{inputenc}
\usepackage[T1]{fontenc}
\usepackage{mathptmx}          % Times New Roman equivalent
\usepackage{setspace}          % line spacing
\usepackage[margin=1in, left=1.5in]{geometry}
\usepackage{graphicx}
\usepackage{booktabs}          % professional tables
\usepackage{longtable}         % multi-page tables
\usepackage{array}
\usepackage{hyperref}
\usepackage{fancyhdr}
\usepackage{titlesec}
\usepackage{caption}
\usepackage{amsmath}
\usepackage{enumitem}
\usepackage{xcolor}
\usepackage{float}

% ── Spacing ─────────────────────────────────────────────────
\onehalfspacing
\setlength{\parindent}{0.5in}
\setlength{\parskip}{6pt}

% ── Chapter / Section Title Formatting ──────────────────────
\titleformat{\chapter}[display]
  {\normalfont\fontsize{16}{20}\bfseries\centering}
  {CHAPTER -- \thechapter}
  {12pt}
  {\fontsize{16}{20}\bfseries\MakeUppercase}

\titleformat{\section}
  {\normalfont\fontsize{14}{16}\bfseries}
  {\thesection}
  {1em}{}

\titleformat{\subsection}
  {\normalfont\fontsize{13}{15}\bfseries}
  {\thesubsection}
  {1em}{}

% ── Header / Footer ─────────────────────────────────────────
\pagestyle{fancy}
\fancyhf{}
\fancyhead[L]{\small\leftmark}
\fancyhead[R]{\small Amrita Vishwa Vidyapeetham}
\fancyfoot[C]{\thepage}
\renewcommand{\headrulewidth}{0.4pt}

% ── Hyperlinks ───────────────────────────────────────────────
\hypersetup{
  colorlinks=true,
  linkcolor=black,
  urlcolor=blue,
  citecolor=black
}

% ════════════════════════════════════════════════════════════
\begin{document}

% ── Title Page ──────────────────────────────────────────────
\begin{titlepage}
  \centering
  \vspace*{0.5cm}

  \includegraphics[width=0.25\textwidth]{amrita_logo.png}\\[0.5cm]
  % Replace amrita_logo.png with the actual logo file

  {\fontsize{14}{18}\bfseries Amrita Vishwa Vidyapeetham}\\[0.2cm]
  {\fontsize{12}{16}\itshape Amrita School of Computing, Coimbatore}\\[0.1cm]
  {\fontsize{12}{16}\itshape Department of Computer Science and Engineering}\\[1cm]

  \rule{\linewidth}{1pt}\\[0.4cm]
  {\fontsize{16}{22}\bfseries
    Secure Blockchain-Based Financial Transaction System\\
    Using Cryptographic Primitives and Smart Contracts}\\[0.4cm]
  \rule{\linewidth}{1pt}\\[1cm]

  {\fontsize{12}{16} \textit{A Project Report submitted in partial fulfillment\\
  of the requirements for the degree of}}\\[0.3cm]
  {\fontsize{13}{17}\bfseries Bachelor of Technology}\\
  {\fontsize{13}{17}\bfseries in Computer Science and Engineering}\\[1cm]

  \begin{tabular}{rl}
    \textbf{Submitted by:}   & Dharaneswara Reddy \\[0.3cm]
    \textbf{Register No:}    & [Your Register Number] \\[0.3cm]
    \textbf{Under the guidance of:} & [Guide Name], [Designation] \\[0.3cm]
  \end{tabular}\\[1.5cm]

  {\fontsize{12}{16} Academic Year: 2025 -- 2026}\\[0.5cm]
  {\fontsize{12}{16} Semester: VI}

  \vfill
\end{titlepage}

% ── Certificate Page ─────────────────────────────────────────
\thispagestyle{empty}
\begin{center}
  {\fontsize{16}{22}\bfseries CERTIFICATE}
\end{center}
\vspace{1cm}

\noindent This is to certify that the project report entitled
\textbf{``Secure Blockchain-Based Financial Transaction System Using Cryptographic Primitives and Smart Contracts''}
submitted by \textbf{Dharaneswara Reddy} (Register No: [Your Register Number])
in partial fulfillment of the requirements for the award of the degree of
\textbf{Bachelor of Technology in Computer Science and Engineering}
at Amrita Vishwa Vidyapeetham, Coimbatore, is a record of bonafide work
carried out by the student under our supervision during the academic year
2025--2026.

\vspace{2cm}

\noindent
\begin{minipage}{0.45\textwidth}
  \centering
  \rule{6cm}{0.4pt}\\
  \textbf{[Guide Name]}\\
  [Designation]\\
  Department of CSE
\end{minipage}
\hfill
\begin{minipage}{0.45\textwidth}
  \centering
  \rule{6cm}{0.4pt}\\
  \textbf{[HOD Name]}\\
  Head of the Department\\
  Department of CSE
\end{minipage}

\vspace{2cm}
\noindent
\begin{center}
  \rule{6cm}{0.4pt}\\
  \textbf{[Principal Name]}\\
  Principal\\
  Amrita School of Computing
\end{center}

\newpage

% ── Abstract ─────────────────────────────────────────────────
\thispagestyle{plain}
\begin{center}
  {\fontsize{16}{22}\bfseries ABSTRACT}
\end{center}
\addcontentsline{toc}{chapter}{ABSTRACT}
\vspace{0.5cm}

\noindent
Financial transactions in centralized banking systems are managed by trusted intermediaries
such as banks and payment processors. These systems are susceptible to fraud, single points
of failure, insider manipulation, and lack of transparency. Blockchain technology offers a
decentralized, tamper-evident alternative where transactions are recorded permanently and
verified without any central authority.

This project presents a secure blockchain-based financial transaction system that implements
the complete transaction lifecycle -- from cryptographic wallet identity creation to on-chain
smart contract settlement. The system implements Elliptic Curve Digital Signature Algorithm
(ECDSA) on the secp256k1 curve for wallet identity and transaction authorization, BIP-62
low-S signature normalization to prevent signature malleability attacks, UUID4-based nonce
binding for replay attack prevention, and SHA-256 Proof-of-Work with dynamic difficulty
adjustment for distributed consensus. Binary Merkle trees are constructed per block to
provide tamper-evident cryptographic commitments to transaction sets. A Solidity smart
contract deployed on an Ethereum-compatible testnet enforces atomic double-spend prevention
using the Checks-Effects-Interactions design pattern.

The complete system is exposed through a RESTful HTTP API and a real-time browser-based
Block Explorer that allows users to submit transactions, trigger mining, and interactively
observe tamper detection. Experimental evaluation confirms that all cryptographic security
invariants hold under adversarial testing, Proof-of-Work convergence matches the
theoretically expected geometric distribution ($16^d$ iterations at difficulty $d$), and
Nakamoto consensus correctly resolves chain forks across multi-node simulation.

\vspace{0.5cm}
\noindent\textbf{Keywords:} Blockchain, ECDSA, secp256k1, Proof-of-Work, Merkle Tree,
Nakamoto Consensus, Smart Contract, Solidity, Replay Attack, Signature Malleability,
Decentralized Finance.

\newpage

% ── Acknowledgements ─────────────────────────────────────────
\thispagestyle{plain}
\begin{center}
  {\fontsize{16}{22}\bfseries ACKNOWLEDGEMENTS}
\end{center}
\addcontentsline{toc}{chapter}{ACKNOWLEDGEMENTS}
\vspace{0.5cm}

\noindent
I would like to express my sincere gratitude to my project guide \textbf{[Guide Name]},
[Designation], Department of Computer Science and Engineering, Amrita Vishwa Vidyapeetham,
for the invaluable guidance, constant encouragement, and constructive feedback throughout
the development of this project. Their deep insights into blockchain technology and
distributed systems greatly enriched the quality and rigor of this work.

I am grateful to the \textbf{Head of the Department}, [HOD Name], for providing the
necessary resources and a conducive research environment. I also thank all the faculty
members of the Department of Computer Science and Engineering for their support and
encouragement during my studies.

I extend my thanks to \textbf{Amrita Vishwa Vidyapeetham} for providing access to the
required computing infrastructure and academic resources that made this project possible.

Finally, I am deeply grateful to my family and friends for their unwavering support,
patience, and motivation throughout the duration of this work.

\vspace{1.5cm}
\begin{flushright}
  \textbf{Dharaneswara Reddy}\\
  B.Tech Computer Science and Engineering\\
  Semester VI\\
  Amrita Vishwa Vidyapeetham, Coimbatore
\end{flushright}

\newpage

% ── Table of Contents ────────────────────────────────────────
\tableofcontents
\newpage

% ── List of Figures ──────────────────────────────────────────
\listoffigures
\newpage

% ── List of Tables ───────────────────────────────────────────
\listoftables
\newpage

% ════════════════════════════════════════════════════════════
% CHAPTER 1 — INTRODUCTION
% ════════════════════════════════════════════════════════════
\chapter{INTRODUCTION}

\section{Introduction to Blockchain Technology}

The concept of blockchain technology was first proposed by Stuart Haber and W. Scott
Stornetta in 1991 as a cryptographically secured chain of blocks to timestamp digital
documents. It became widely known in 2008 when Satoshi Nakamoto published the Bitcoin
whitepaper, introducing a peer-to-peer electronic cash system that eliminated the need
for a trusted financial intermediary \cite{nakamoto2008}. Since then, blockchain
technology has evolved significantly and is now being applied in domains far beyond
cryptocurrency -- including supply chain management, healthcare records, digital identity
verification, insurance claim processing, and financial settlements \cite{zheng2017}.

A blockchain is a distributed, append-only data structure consisting of a sequence of
cryptographically linked blocks. Each block contains a set of transactions, a timestamp,
a nonce value used for mining, and the SHA-256 hash of the preceding block. This backward
hash pointer creates a dependency chain: modifying any historical block changes its hash,
which invalidates the hash pointer stored in the next block, and this invalidation cascades
through all subsequent blocks, making tampering immediately detectable.

The fundamental innovation of blockchain is its ability to achieve distributed consensus
among mutually distrusting participants without relying on any central authority. This
property is enforced through the Proof-of-Work (PoW) consensus mechanism, in which nodes
compete to find a nonce such that the hash of a block satisfies a computational difficulty
condition. This mechanism makes it economically infeasible for any single party to rewrite
the history of the ledger without controlling a majority of the network's total computing
power.

\section{Motivation}

The traditional financial system depends heavily on centralized intermediaries -- banks,
clearinghouses, and payment processors -- to authorize and record transactions. While
these systems are well-established, they suffer from three structural vulnerabilities.

First, they represent a single point of failure. If the central authority is compromised,
corrupted, or simply unavailable, every account and transaction it manages is at risk.
High-profile breaches of centralized financial institutions have demonstrated that this
risk is real and consequential.

Second, users of centralized systems have no independent means of verifying that their
account balances are accurate or that their transaction history has not been manipulated.
The entire trust model rests on institutional integrity rather than cryptographic proof.

Third, the cost and latency of cross-border financial settlement remain high due to the
number of intermediaries that must coordinate: a typical international wire transfer
passes through several correspondent banks and clearinghouse systems before reaching its
destination, taking two to five business days and incurring significant fees.

Blockchain technology addresses all three of these structural problems simultaneously. The
distributed ledger eliminates the single point of failure. Cryptographic hash chaining and
digital signatures replace institutional trust with mathematical proof. And the near-instant
block finality of local blockchain networks replaces slow multi-party settlement with
transparent, immediate confirmation.

However, the practical adoption of blockchain financial systems requires engineers who
understand not merely the high-level architecture but the specific security controls that
make blockchain transactions secure. The most widely cited blockchain educational
implementations consistently omit these controls, creating a gap between conceptual
understanding and secure practice. This project is designed to bridge that gap.

\section{Objectives}

The primary objectives of this project are:

\begin{enumerate}[label=\arabic*.]
  \item To design and implement a complete blockchain-based financial transaction system
        that operates without any trusted central authority.
  \item To implement ECDSA digital signatures on the secp256k1 curve for wallet identity
        and transaction authorization, with BIP-62 low-S normalization to prevent signature
        malleability attacks.
  \item To prevent replay attacks through UUID4-based nonce binding embedded directly in
        the signed transaction payload.
  \item To implement SHA-256 Proof-of-Work consensus with dynamic difficulty adjustment
        for stable block production.
  \item To construct binary Merkle trees per block, providing cryptographic commitment to
        all transactions and enabling efficient tamper detection.
  \item To deploy a Solidity smart contract on an Ethereum-compatible testnet that enforces
        atomic balance management and prevents double-spending through the
        Checks-Effects-Interactions design pattern.
  \item To expose the system through a RESTful API and a real-time browser-based Block
        Explorer that enables live demonstration of mining, consensus, and tamper detection.
  \item To enforce software quality through an automated CI/CD pipeline with static security
        analysis, CVE dependency auditing, and minimum 80\% test coverage.
\end{enumerate}

\section{Scope of the Project}

This project implements a fully functional, self-contained blockchain financial transaction
network at an educational simulation scale. The scope includes:

\begin{itemize}
  \item Cryptographic wallet identity management using real Ethereum addresses derived from
        a public transaction dataset.
  \item A complete transaction signing and verification protocol with four independent
        security controls (ECDSA, BIP-62, UUID4 nonce, and canonical JSON serialization).
  \item A blockchain engine with SHA-256 PoW mining, Merkle tree construction, mempool
        management, and O($n$) chain integrity validation.
  \item A simulated peer-to-peer network with Nakamoto longest-chain consensus and
        fork resolution.
  \item A Solidity smart contract implementing wallet funding, transaction submission,
        approval, and rejection with reentrancy protection.
  \item A REST API with nine endpoints and a browser-based Block Explorer dashboard.
\end{itemize}

The system is designed for local deployment using a Ganache testnet and is not intended
for production mainnet use without additional infrastructure hardening.

\section{Organization of the Report}

The remainder of this report is organized as follows. Chapter 2 presents a review of
related literature. Chapter 3 describes the hardware and software requirements. Chapter 4
presents the system design, including high-level and low-level architecture. Chapter 5
describes the system implementation. Chapter 6 presents system testing results. Chapter 7
provides results and analysis. Chapter 8 concludes the report and outlines future work.

% ════════════════════════════════════════════════════════════
% CHAPTER 2 — LITERATURE SURVEY
% ════════════════════════════════════════════════════════════
\chapter{LITERATURE SURVEY}

This chapter reviews the most significant prior work in the domains of blockchain
architecture, cryptographic digital signatures, Proof-of-Work consensus, smart contract
security, and decentralized financial systems. The survey identifies the open problems
that motivate the design decisions made in this project.

\section{Foundational Blockchain Architecture}

Nakamoto \cite{nakamoto2008} introduced the first complete decentralized payment system
in 2008. The Bitcoin protocol demonstrated that digital cash could be transferred between
parties without a trusted intermediary by combining SHA-256 Proof-of-Work for Sybil
resistance, a chain of cryptographic hash pointers for tamper evidence, and the
longest-chain rule for distributed consensus. The UTXO (Unspent Transaction Output) model
was used to track balances without maintaining account state. While Bitcoin established
the fundamental architecture that all subsequent blockchain systems build upon, its
scripting language is intentionally limited and cannot express complex financial logic
such as conditional payments, multi-party fund management, or programmable settlement rules.

Zheng et al. \cite{zheng2017} conducted a comprehensive survey of blockchain architectures,
consensus mechanisms, and application domains across the rapidly growing literature.
Their analysis categorized consensus algorithms into Proof-of-Work, Proof-of-Stake, and
Byzantine fault tolerant (PBFT) families and quantified the tradeoffs in terms of
throughput, finality, and security guarantees. Their survey confirms that Proof-of-Work
remains the most battle-tested mechanism for achieving trustless consensus in a
permissionless network, despite its energy cost, because it does not require a pre-selected
validator set and is therefore not vulnerable to Sybil attacks in the same way that
Proof-of-Stake systems can be.

\section{Elliptic Curve Digital Signatures and Malleability}

Johnson, Menezes, and Vanstone \cite{johnson2001} provided the formal mathematical
specification and security proofs for ECDSA. They demonstrated that ECDSA signatures
are existentially unforgeable against chosen-message attacks under the random oracle model,
given the hardness of the elliptic curve discrete logarithm problem. The secp256k1 curve,
selected for both Bitcoin and Ethereum due to its transparent parameter generation and
well-studied security properties, operates over a 256-bit prime field with a generator
point of prime order $N = 2^{256} - 432420386565659656852420866394968145599$.

A critical implementation vulnerability of ECDSA -- signature malleability -- was
formally documented in Bitcoin Improvement Proposal BIP-62 \cite{bip62}. Because the
ECDSA mathematical structure produces two valid signatures for every signing operation
($(r, s)$ and $(r, N-s)$), an adversary who observes a valid signed transaction can
produce an alternate valid signature without possessing the private key. Since the
transaction identifier was historically derived from the serialized transaction including
its signature, this allowed an attacker to modify the transaction ID without invalidating
the authorization -- a property exploited in the MtGox exchange collapse of 2014,
resulting in losses of approximately 460 million USD. BIP-62 proposed low-S normalization
as the canonical fix: signatures must satisfy $s \leq N/2$, with $N-s$ substituted when
this condition is violated. This standard was subsequently adopted by Ethereum as well.

\section{Proof-of-Work Security and Difficulty Adjustment}

Gervais et al. \cite{gervais2016} performed a systematic quantitative security analysis
of Proof-of-Work blockchains. Their work modeled blockchain systems as stochastic
processes and quantified the relationship between block production rate, network
propagation delay, and resistance to double-spend and selfish mining attacks. A key
finding was that overly fast block production (relative to propagation delay) increases
the rate of honest block orphaning, which reduces effective chain security even without
an active attacker. This motivates the use of a dynamic difficulty adjustment mechanism
to maintain a stable block production rate regardless of available hash power.

Eyal and Sirer \cite{eyal2014} demonstrated that the Nakamoto honest majority assumption
is weaker than commonly believed. Through the selfish mining strategy, a pool controlling
as little as 25\% to 33\% of total hash rate can gain a disproportionate share of mining
rewards by withholding discovered blocks strategically. This finding underscores the
importance of difficulty adjustment mechanisms that prevent mining pools from gaining
sustained rate advantages and identifies the 51\% attack as a genuine boundary condition
for single-node simulations.

\section{Smart Contract Security}

Buterin \cite{buterin2014} introduced Ethereum as a programmable extension of the
blockchain model. By introducing the Ethereum Virtual Machine (EVM) and the Solidity
programming language, Ethereum enabled developers to write smart contracts -- programs
that execute on the blockchain deterministically across all nodes and can implement
arbitrary financial rules without a central operator. This created the foundation for
decentralized finance (DeFi), where lending, trading, and settlement occur without banks.

Luu et al. \cite{luu2016} identified and formally demonstrated reentrancy as the most
dangerous vulnerability class in Solidity smart contracts. In a reentrancy attack, a
malicious contract recursively calls a victim function before its first invocation
completes, exploiting the window between the balance check and the balance update to drain
funds multiple times from a single authorization. This vulnerability was exploited in the
2016 DAO hack, resulting in the loss of approximately 60 million USD. Their work
established the Checks-Effects-Interactions (CEI) pattern as the standard defense:
all balance deductions (Effects) must occur before any external calls or event emissions
(Interactions), so that any re-entrant call fails the balance check (Checks) immediately.

\section{Double-Spending Prevention}

The 2021 IEEE World AI IoT Congress \cite{ieee2021bar} presented the Block Access
Restriction (BAR) mechanism to mitigate 51\% attacks and double-spending in PoW networks
by monitoring miner behavior and restricting block access privileges upon detection of
anomalous activity. The authors of \cite{ieee2021sybil} provided a security risk
management framework for systematically analyzing Sybil and double-spending attack surfaces
in blockchain networks, categorizing attack vectors and proposed countermeasures.
The work of \cite{mdpi2021hybrid} proposed combining PoW and Proof-of-Stake consensus
mechanisms to increase the computational and economic cost of double-spend attacks beyond
what pure PoW systems can achieve.

\section{Summary and Research Gap}

Table \ref{tab:litsurvey} summarizes the literature reviewed and identifies the open
problem addressed by each reference in the context of this project.

\begin{longtable}{|p{0.4cm}|p{3.0cm}|p{5.0cm}|p{3.5cm}|}
\caption{Literature Survey Summary}
\label{tab:litsurvey}\\
\hline
\textbf{S.} & \textbf{Author(s)} & \textbf{Title and Year} & \textbf{Open Problem / Inference} \\
\hline
\endfirsthead
\hline
\textbf{S.} & \textbf{Author(s)} & \textbf{Title and Year} & \textbf{Open Problem / Inference} \\
\hline
\endhead
1 & S. Nakamoto & Bitcoin: A Peer-to-Peer Electronic Cash System (2008) & Limited scripting; no smart contracts or complex financial logic \\
\hline
2 & Z. Zheng et al. & An Overview of Blockchain Technology: Architecture, Consensus, and Future Trends -- IEEE BigData (2017) & PoW is most secure but energy-intensive; no unified implementation covering all security controls \\
\hline
3 & D. Johnson et al. & The Elliptic Curve Digital Signature Algorithm (ECDSA) -- IJIS (2001) & Standard ECDSA is vulnerable to signature malleability; BIP-62 normalization must be applied manually \\
\hline
4 & V. Buterin & Ethereum: A Next-Generation Smart Contract Platform (2014) & Smart contracts vulnerable to reentrancy; CEI pattern not automatically enforced \\
\hline
5 & A. Gervais et al. & On the Security and Performance of Proof of Work Blockchains -- ACM CCS (2016) & Dynamic difficulty adjustment needed to maintain stable block rate and security \\
\hline
6 & L. Luu et al. & Making Smart Contracts Smarter -- ACM CCS (2016) & Reentrancy allows fund drainage; balance deduction must precede any state write \\
\hline
7 & I. Eyal, E. G. Sirer & Majority is Not Enough: Bitcoin Mining is Vulnerable -- FC (2014) & Selfish mining effective at 25--33\% hash rate; honest majority assumption is weaker than stated \\
\hline
8 & IEEE AIIoT (2021) & A Protecting Mechanism Against Double Spending Attack in Blockchain Systems (2021) & Single countermeasure insufficient; protocol and contract-level controls both required \\
\hline
9 & IEEE (2021) & Exploring Sybil and Double-Spending Risks in Blockchain Systems (2021) & Frameworks lack integration with smart contract settlement for end-to-end prevention \\
\hline
10 & MDPI Future Internet (2021) & Distributed Hybrid Double-Spending Attack Prevention for PoW and PoS Consensuses (2021) & Hybrid consensus more secure but complex; PoW with atomic smart contract settlement is practical \\
\hline
\end{longtable}

The key finding from the literature survey is that no existing educational or reference
implementation simultaneously addresses all six security vulnerabilities identified --
transaction forgery, replay attacks, signature malleability, hash non-determinism, smart
contract reentrancy, and undetectable chain tampering -- within a single tested and
automated system. This project fills that gap.

% ════════════════════════════════════════════════════════════
% CHAPTER 3 — SYSTEM REQUIREMENTS AND ANALYSIS
% ════════════════════════════════════════════════════════════
\chapter{SYSTEM REQUIREMENTS AND ANALYSIS}

\section{Software Requirements}

The system is built entirely from open-source components. The software requirements are
categorized by layer.

\subsection{Programming Languages and Runtimes}

\begin{table}[H]
\centering
\caption{Programming Languages and Runtime Environments}
\label{tab:languages}
\begin{tabular}{|l|l|p{7cm}|}
\hline
\textbf{Language / Runtime} & \textbf{Version} & \textbf{Purpose} \\
\hline
Python & 3.13.0 & Core blockchain engine, REST API, cryptographic signing, key management \\
\hline
Solidity & 0.8.20 & Smart contract implementation for on-chain settlement \\
\hline
JavaScript & ES2022 & Hardhat deployment scripts; Block Explorer frontend \\
\hline
Node.js & 22 LTS & Hardhat compilation and deployment runtime \\
\hline
\end{tabular}
\end{table}

\subsection{Python Libraries -- Production Dependencies}

\begin{table}[H]
\centering
\caption{Python Production Dependencies}
\label{tab:pydeps}
\begin{tabular}{|l|l|p{7cm}|}
\hline
\textbf{Library} & \textbf{Version} & \textbf{Purpose} \\
\hline
\texttt{cryptography} & $\geq$ 44.0.2 & ECDSA secp256k1, AES-256-CBC encryption, PBKDF2-HMAC key derivation \\
\hline
\texttt{flask} & $\geq$ 3.1.0 & REST API framework; application factory pattern \\
\hline
\texttt{web3} & $\geq$ 7.14.1 & Ethereum JSON-RPC client for smart contract interaction \\
\hline
\texttt{pandas} & $\geq$ 2.3.3 & CSV dataset loading and Ethereum address extraction \\
\hline
\texttt{python-dotenv} & $\geq$ 1.1.0 & Secure environment variable management via .env file \\
\hline
\end{tabular}
\end{table}

\subsection{Blockchain and Smart Contract Tools}

\begin{table}[H]
\centering
\caption{Smart Contract Development Tools}
\label{tab:contracttools}
\begin{tabular}{|l|l|p{7cm}|}
\hline
\textbf{Tool} & \textbf{Version} & \textbf{Purpose} \\
\hline
Hardhat & $\wedge$2.22.18 & Solidity compilation, local testnet management, deployment \\
\hline
Ganache GUI & 2.7.1 & Local Ethereum testnet (Chain ID 1337, port 7545) \\
\hline
ethers.js & $\wedge$6.13.4 & Ethereum client library for deployment scripts \\
\hline
\end{tabular}
\end{table}

\subsection{Development and Quality Assurance Tools}

\begin{table}[H]
\centering
\caption{Development and CI/CD Tools}
\label{tab:devtools}
\begin{tabular}{|l|p{9cm}|}
\hline
\textbf{Tool} & \textbf{Purpose} \\
\hline
pytest $+$ pytest-cov & Unit and integration test runner with statement coverage measurement \\
\hline
Bandit & Static security analysis for Python source code; detects injection, hardcoded secrets, and unsafe calls \\
\hline
pip-audit & Dependency CVE scanning against the National Vulnerability Database \\
\hline
Ruff & Python linter and code formatter enforcing PEP 8 and security rules \\
\hline
GitHub Actions & Automated five-stage CI/CD pipeline triggered on every commit to main branch \\
\hline
\end{tabular}
\end{table}

\section{Hardware Requirements}

The system is designed to run on standard commodity hardware. No GPU or specialized
cryptographic accelerator is required.

\begin{table}[H]
\centering
\caption{Hardware Requirements}
\label{tab:hardware}
\begin{tabular}{|l|l|l|}
\hline
\textbf{Component} & \textbf{Minimum} & \textbf{Recommended} \\
\hline
Operating System & Ubuntu 20.04 / Windows 10 / macOS 12 & Ubuntu 22.04 LTS \\
\hline
Processor & Any modern x86-64 CPU & Intel Core i5 / AMD Ryzen 5 or better \\
\hline
RAM & 4 GB & 8 GB \\
\hline
Storage & 500 MB free & 2 GB free (for blockchain state, logs) \\
\hline
Network & Localhost loopback only & 100 Mbps (for testnet access) \\
\hline
\end{tabular}
\end{table}

\section{Problem Analysis}

The analysis of the problem domain identifies six specific technical attack vectors
that must be addressed in any secure blockchain financial transaction system. These
are formally defined as follows.

\textbf{P1 -- Transaction Forgery:} Without cryptographic signing, any network participant
can submit a transaction claiming to spend funds owned by another party. The system must
provide a mechanism by which the spending authority of any transaction can be
cryptographically verified by any node without trusting the sender.

\textbf{P2 -- Replay Attacks:} A validly signed transaction is a fixed byte sequence.
Once broadcast, it can be captured and resubmitted by an adversary at any time. Without
per-transaction uniqueness embedded in the signature itself, a single authorization can
trigger multiple executions of the same transfer.

\textbf{P3 -- Signature Malleability:} The ECDSA mathematical structure produces two
valid signatures for every signing operation: $(r, s)$ and $(r, N-s)$. An adversary who
intercepts a signed transaction can substitute the alternate signature and present it as
a different transaction with a different identifier, enabling fraudulent claims of
non-execution.

\textbf{P4 -- Hash Non-Determinism:} Block and transaction hashing requires byte-level
determinism across all environments. Languages that do not guarantee deterministic
serialization of data structures (such as unordered dictionary iteration in Python) will
produce different hash values for the same logical data in different execution contexts,
silently corrupting the integrity guarantees of the hash chain.

\textbf{P5 -- Smart Contract Reentrancy and Double-Spending:} A smart contract that
checks a sender's balance, then initiates a transfer, then updates the balance creates
a window during which a malicious caller can recursively re-enter the function. Each
recursive call passes the balance check (because the balance has not yet been updated)
and transfers funds again, resulting in multiple debits from a single authorization.

\textbf{P6 -- Undetectable Chain Tampering:} The system must provide a formally
verifiable, automatic mechanism to detect any post-hoc modification of any historical
transaction without relying on external audit or trusted witnesses. Detection must
propagate immediately through the entire chain.

% ════════════════════════════════════════════════════════════
% CHAPTER 4 — SYSTEM DESIGN
% ════════════════════════════════════════════════════════════
\chapter{SYSTEM DESIGN}

\section{High-Level Design}

The proposed system is organized as a six-layer architecture in which each layer encapsulates
a single responsibility and communicates with adjacent layers through well-defined interfaces.
This strict layering ensures that each security control can be tested and verified
independently. Figure \ref{fig:arch} shows the high-level system architecture.

\begin{figure}[H]
  \centering
  % Insert diagram_ppt.png here
  \fbox{\parbox{0.92\textwidth}{\centering\vspace{2cm}
  [Insert System Architecture Diagram -- diagram\_ppt.png]\vspace{2cm}}}
  \caption{Six-Layer System Architecture of the Blockchain Financial Transaction System}
  \label{fig:arch}
\end{figure}

The six layers, from bottom to top, are:

\begin{enumerate}[label=\arabic*.]
  \item \textbf{Cryptographic Identity Layer} -- ECDSA secp256k1 key pair generation,
        AES-256-CBC encrypted private key storage, and Ethereum-format address validation.
  \item \textbf{Transaction Signing Layer} -- ECDSA-SHA256 signing with BIP-62 low-S
        normalization, UUID4 nonce binding, and canonical JSON serialization.
  \item \textbf{Blockchain Engine Layer} -- SHA-256 Proof-of-Work mining, binary Merkle
        tree construction, FIFO mempool management, and chain integrity validation.
  \item \textbf{Peer-to-Peer Consensus Layer} -- Nakamoto longest-chain rule, block
        broadcast, and fork resolution across multiple nodes.
  \item \textbf{Smart Contract Layer} -- Solidity contract on Ethereum-compatible testnet;
        wallet funding, transaction submission, approval, and rejection with CEI pattern.
  \item \textbf{Presentation Layer} -- Flask REST API with nine endpoints and a browser-based
        Block Explorer with real-time chain monitoring.
\end{enumerate}

\section{Low-Level Design}

\subsection{Cryptographic Identity Design}

The identity of every participant in the financial network is established by an ECDSA
key pair generated using the secp256k1 elliptic curve. The curve is defined over a 256-bit
prime field $\mathbb{F}_p$ where:
\[
  p = 2^{256} - 2^{32} - 977
\]
The generator point $G$ has prime order $N \approx 2^{256}$. A private key is a randomly
sampled scalar $k \in \mathbb{Z}_N$ drawn from a cryptographically secure random number
generator backed by the operating system entropy pool. The corresponding public key is the
curve point $Q = k \cdot G$, computed via elliptic curve scalar multiplication.

The wallet address is derived following the Ethereum convention: the uncompressed public
key (64 bytes: $x || y$ coordinates) is hashed with Keccak-256, and the last 20 bytes of
the result are taken as the address. This produces a 40-character hexadecimal string
prefixed with \texttt{0x}, which is the Ethereum address format.

Private keys are stored encrypted at rest using AES-256-CBC symmetric encryption, with
the encryption key derived from a user-supplied passphrase through PBKDF2-HMAC-SHA256 key
derivation (100,000 iterations). The encrypted key material is stored in PKCS8 PEM format
with filesystem permissions set to \texttt{0o600}, ensuring that the file is readable only
by its owner. The passphrase is never stored on disk or written to log files.

\subsection{Transaction Signing Protocol}

The transaction signing protocol implements four security controls in sequence:

\textbf{Step 1 -- Nonce Generation (P2):} A UUID4 value is generated using 122 bits of
OS entropy via \texttt{os.urandom}. This nonce is embedded directly in the transaction
data structure before signing. Because the nonce is part of the signed payload, any attempt
to reuse the signature with a different nonce produces a different message and therefore
fails verification.

\textbf{Step 2 -- Canonical Serialization (P4):} The transaction data (including the
embedded nonce) is serialized to JSON with alphabetically sorted keys and compact separators
(\texttt{sort\_keys=True}, \texttt{separators=(',',':')}). This canonical form guarantees
byte-level determinism: the same logical transaction always produces the same byte sequence
regardless of the programming environment, dictionary insertion order, or platform.

\textbf{Step 3 -- ECDSA Signing (P1):} The canonical JSON bytes are signed using
ECDSA-SHA256 with the sender's private key. The implementation uses RFC 6979 deterministic
nonce generation for the per-signature ephemeral key $k$, eliminating the risk of private
key recovery through weak randomness.

\textbf{Step 4 -- Low-S Normalization (P3):} After signing, if the $s$ component of the
signature satisfies $s > N/2$, it is replaced with $N - s$. The resulting signature is
serialized as a fixed 64-byte value: $r$ (32 bytes, big-endian) concatenated with $s$
(32 bytes, big-endian). The verifier enforces the complementary condition, rejecting any
signature with $s > N/2$ with a \texttt{VerificationError}. This makes both members of
the ECDSA mathematical twin pair resolve to the same canonical form.

\subsection{Blockchain Data Structure}

Each block in the chain is a data record containing the following fields:

\begin{table}[H]
\centering
\caption{Block Data Structure}
\label{tab:blockstruct}
\begin{tabular}{|l|l|p{7cm}|}
\hline
\textbf{Field} & \textbf{Type} & \textbf{Description} \\
\hline
index & Integer & Zero-based position of the block in the chain \\
\hline
transactions & List & Ordered list of financial transaction records \\
\hline
merkle\_root & String & Root of the binary Merkle tree over all transactions \\
\hline
previous\_hash & String & SHA-256 hash of the preceding block (chain link) \\
\hline
nonce & Integer & Value found by PoW miner satisfying difficulty condition \\
\hline
timestamp & Float & Unix timestamp of block creation \\
\hline
hash & String & SHA-256 hash of all above fields (canonical JSON) \\
\hline
\end{tabular}
\end{table}

The genesis block (index 0) has \texttt{previous\_hash} set to a string of 64 zeros,
establishing the base case of the chain. Every subsequent block's \texttt{previous\_hash}
must exactly equal the \texttt{hash} of the preceding block, creating the chain dependency.

\subsection{Merkle Tree Construction}

Figure \ref{fig:merkle} illustrates the Merkle tree construction for a block containing
three transactions $T_0, T_1, T_2$.

\begin{figure}[H]
  \centering
  \fbox{\parbox{0.88\textwidth}{\centering\vspace{2cm}
  [Insert Merkle Tree Diagram -- Figure showing leaf hashing and pairwise combination]\vspace{2cm}}}
  \caption{Binary Merkle Tree Construction for a 3-Transaction Block}
  \label{fig:merkle}
\end{figure}

The algorithm proceeds as follows:
\begin{enumerate}[label=\arabic*.]
  \item Each transaction $T_i$ is serialized to canonical JSON and hashed with SHA-256
        to produce a leaf node $L_i = \text{SHA256}(T_i)$.
  \item If the number of leaves is odd, the last leaf is duplicated so that each level
        has an even number of nodes.
  \item Pairs of adjacent nodes are concatenated and hashed: $P_{ij} = \text{SHA256}(L_i \| L_j)$.
  \item This process repeats up the tree until a single root value $R$ remains.
  \item The root $R$ is stored in the block and incorporated into the block hash.
\end{enumerate}

Any modification to any transaction changes its leaf hash, which cascades up the tree to
change the root, which changes the block hash, which breaks the \texttt{previous\_hash}
pointer in the next block, making tamper detection automatic and O($n$) over the chain.

\subsection{Proof-of-Work Algorithm}

Mining is the process of finding a nonce $n$ such that:
\[
  \text{SHA256}\bigl(\text{canonical\_JSON}(\text{block with nonce } n)\bigr)
  \text{ starts with } \underbrace{00\ldots0}_{d \text{ zeros}}
\]
where $d$ is the current network difficulty parameter. The probability of any single
hash evaluation satisfying this condition is $p = 16^{-d}$, making the expected number
of iterations $E[\text{iters}] = 16^d$. At $d = 3$, the expected iterations are 4,096.

The dynamic difficulty adjustment rule is evaluated every 10 blocks:
\[
  \text{new difficulty} = \begin{cases}
    d + 1 & \text{if elapsed time for last 10 blocks} < 50\text{s (too fast)} \\
    d - 1 & \text{if elapsed time for last 10 blocks} > 200\text{s (too slow)} \\
    d     & \text{otherwise}
  \end{cases}
\]
A minimum difficulty of 1 is enforced to prevent the chain from stalling. This mirrors
Bitcoin's retargeting mechanism in simplified form, maintaining a stable block rate as
available hash power changes.

\subsection{Nakamoto Consensus and Fork Resolution}

Each network node maintains a list of peer nodes. When a node successfully mines a block,
it broadcasts the block to all peers via direct method invocation (simulated P2P). Each
receiving node validates two conditions:
\begin{enumerate}[label=\arabic*.]
  \item The block's hash satisfies the current PoW difficulty condition.
  \item The block's \texttt{previous\_hash} equals the hash of the receiving node's
        current chain tip.
\end{enumerate}

If both conditions hold, the block is appended. If they fail because the peer's chain is
longer, the receiving node initiates chain synchronization using the Nakamoto rule:

\[
  \text{Replace local chain if:} \quad
  |\text{peer.chain}| > |\text{local.chain}| \;\wedge\; \text{peer.is\_valid\_chain()} = \text{True}
\]

The double condition (longer \textit{and} valid) prevents an adversary from forcing a
chain replacement by broadcasting a long but structurally invalid chain.

\subsection{Smart Contract Design}

The Solidity smart contract implements on-chain financial settlement with four functions,
as shown in Figure \ref{fig:contract}.

\begin{figure}[H]
  \centering
  \fbox{\parbox{0.88\textwidth}{\centering\vspace{2cm}
  [Insert Smart Contract State Machine Diagram]\vspace{2cm}}}
  \caption{Smart Contract Transaction State Machine}
  \label{fig:contract}
\end{figure}

The critical security design is the application of the Checks-Effects-Interactions (CEI)
pattern in the \texttt{submitTransaction} function:

\begin{itemize}
  \item \textbf{CHECKS:} Verify that \texttt{balances[sender] >= value + fee}.
  \item \textbf{EFFECTS:} Deduct \texttt{balances[sender] -= (value + fee)} \textit{before}
        any other operation.
  \item \textbf{INTERACTIONS:} Write the transaction record and emit the confirmation event.
\end{itemize}

This ordering closes the reentrancy window (P5): any recursive call to the function will
find the sender's balance already reduced and fail the balance check immediately.

\subsection{REST API Design}

\begin{table}[H]
\centering
\caption{REST API Endpoint Reference}
\label{tab:api}
\begin{tabular}{|l|l|p{6cm}|}
\hline
\textbf{Method} & \textbf{Endpoint} & \textbf{Description} \\
\hline
POST & /api/transactions/submit & Add a signed transaction to the mempool \\
\hline
GET  & /api/transactions/pending & Retrieve all unconfirmed transactions \\
\hline
POST & /api/mine & Execute PoW mining on pending transactions \\
\hline
GET  & /api/blocks & Retrieve the full blockchain \\
\hline
GET  & /api/blocks/\{index\} & Retrieve a single block by index \\
\hline
GET  & /api/chain/stats & Chain height, difficulty, validity, pending count \\
\hline
GET  & /api/chain/validate & Full chain integrity check \\
\hline
POST & /api/tamper/\{index\} & Inject fake transaction for demo (P6 demonstration) \\
\hline
POST & /api/chain/restore & Reset chain to clean genesis state \\
\hline
\end{tabular}
\end{table}

% ════════════════════════════════════════════════════════════
% CHAPTER 5 — SYSTEM IMPLEMENTATION
% ════════════════════════════════════════════════════════════
\chapter{SYSTEM IMPLEMENTATION}

\section{Module 1: Dataset and Address Loading}

The wallet identity pool is initialized from a real Ethereum mainnet transaction dataset
(\texttt{Eth\_Txs.csv}) containing historical sender and receiver addresses from public
blockchain transactions. The address extraction pipeline processes this dataset through
five sequential stages.

\textbf{Stage 1 -- File Validation:} The CSV file is checked for existence and readability.
A typed \texttt{DatasetLoadError} exception is raised if the file cannot be opened, providing
a clear fail-fast signal.

\textbf{Stage 2 -- Column Selection:} Only the \texttt{From} and \texttt{To} columns are
loaded, with all values forced to the string type using \texttt{dtype=str} to prevent
pandas from misinterpreting hexadecimal address strings as numeric values.

\textbf{Stage 3 -- Address Extraction:} A set union is computed over all non-null values
in both columns: \texttt{set(df["From"].dropna()) | set(df["To"].dropna())}. Set semantics
automatically deduplicate the addresses.

\textbf{Stage 4 -- Regex Validation:} Each candidate address is matched against the
compiled pattern \texttt{\^{}0x[0-9a-fA-F]\{40\}\$}. Addresses that do not match
(exchange labels, ENS names, and other non-address values) are logged at \texttt{WARNING}
level and excluded. Valid addresses are retained.

\textbf{Stage 5 -- Normalization:} All valid addresses are converted to lowercase for
consistent comparison throughout the system, as Ethereum address comparisons are
case-insensitive.

The output is a \texttt{set[str]} of unique, validated, lowercase Ethereum addresses used
as the universe of identities for key pair generation.

\section{Module 2: Key Management}

For each wallet address in the loaded set, an ECDSA key pair is generated as follows.
A 256-bit private key scalar is sampled from the OS cryptographic random number generator.
The corresponding public key is computed via elliptic curve scalar multiplication on the
secp256k1 curve.

The private key is encrypted using AES-256-CBC with a key derived from the configured
passphrase via PBKDF2-HMAC-SHA256 (100,000 iterations, 32-byte output). The encrypted
key material is serialized to PKCS8 PEM format and written to the keystore directory.
The file permission is immediately set to \texttt{0o600} using \texttt{os.chmod}, ensuring
that no other user account on the system can read the key material.

Loading a key requires the same passphrase. If an incorrect passphrase is provided, the
decryption fails with a \texttt{ValueError} from the cryptography library, which is caught
and re-raised as a typed \texttt{KeyStorageError}. This ensures that authentication
failures produce descriptive, catchable exceptions rather than cryptic low-level errors.

\section{Module 3: Transaction Signing and Verification}

\subsection{Signing}

The signing function accepts a transaction dictionary, the sender's private key object,
and a flag indicating whether to generate a new nonce or use a caller-supplied one. The
implementation:

\begin{enumerate}[label=\arabic*.]
  \item Calls \texttt{create\_nonce()} to generate a UUID4 string from \texttt{os.urandom(16)}.
  \item Inserts the nonce into the transaction dictionary under the key \texttt{"\_\_nonce\_\_"}.
  \item Serializes the augmented dictionary to bytes using canonical JSON
        (\texttt{sort\_keys=True}, \texttt{separators=(',',':')}).
  \item Signs the bytes using \texttt{private\_key.sign(message, ECDSA(), SHA256())}.
  \item Decodes the DER-encoded signature to extract the raw integer pair $(r, s)$.
  \item Applies BIP-62 normalization: if $s > N/2$, sets $s \leftarrow N - s$.
  \item Re-encodes to the fixed 64-byte format: $r$ as 32 bytes big-endian followed by
        $s$ as 32 bytes big-endian.
  \item Returns the signature bytes and the nonce string.
\end{enumerate}

\subsection{Verification}

The verification function accepts a signature, public key, transaction dictionary, and nonce.
It enforces four sequential checks:

\begin{enumerate}[label=\arabic*.]
  \item Length check: \texttt{len(sig) == 64}. Rejects any non-standard encoding.
  \item BIP-62 check: extracts $s = \text{int.from\_bytes}(\text{sig[32:]}, \text{'big'})$
        and raises \texttt{VerificationError} if $s > N/2$.
  \item Canonical reconstruction: rebuilds the message bytes from the transaction and
        nonce using the same canonical JSON serialization as signing.
  \item ECDSA verification: re-encodes $(r, s)$ to DER format and calls
        \texttt{public\_key.verify()}, raising \texttt{VerificationError} on
        \texttt{InvalidSignature}.
\end{enumerate}

\section{Module 4: Blockchain Engine}

\subsection{Block Creation and Hashing}

A block is instantiated as a Python dataclass. The \texttt{\_\_post\_init\_\_} method
immediately computes the Merkle root from the transaction list and the block hash from all
fields using canonical JSON serialization, ensuring that the hash is always consistent
from the moment of creation and cannot be accidentally out of sync with the block content.

\subsection{Mempool}

The mempool is implemented using \texttt{collections.deque} with a configured maximum
length of 1000 transactions. Each incoming transaction is stamped with a UUID4 identifier
and a Unix timestamp on entry. The deque provides O(1) append and O(1) left-pop (FIFO
selection), ensuring that transaction submission and mining selection are both constant-time
operations regardless of mempool size.

\subsection{Mining}

The mining function selects up to 10 transactions from the mempool using FIFO order and
prepends a coinbase transaction crediting the mining reward to the miner's address. It
then enters a nonce-search loop, incrementing the nonce from zero, computing the candidate
block hash at each step, and testing the leading-zero prefix condition. When a valid nonce
is found, the block is passed to the blockchain's \texttt{add\_block()} method.

\subsection{Chain Validation (P6)}

The \texttt{is\_valid\_chain()} function iterates through the chain from index 1 and checks
two invariants for every block:
\begin{enumerate}[label=\arabic*.]
  \item \textbf{Stored hash consistency:} \texttt{block.hash == block.compute\_hash()}. This
        detects any mutation of block content after mining.
  \item \textbf{Chain linkage:} \texttt{block.previous\_hash == chain[i-1].hash}. This detects
        any break in the backward pointer chain.
\end{enumerate}
If either check fails for any block, the function returns \texttt{False} immediately. This
implements P6 resolution with O($n$) time complexity over the chain length $n$.

\section{Module 5: Smart Contract}

The Solidity smart contract (\texttt{TransactionContract.sol}) manages an internal ledger
of user balances as a \texttt{mapping(address => uint256)}. The \texttt{fundWallet()}
function accepts ETH and credits the internal balance for any specified address.

The \texttt{submitTransaction()} function is the most security-critical. It first checks
that the caller's internal balance is sufficient for the value plus the fee. It then
immediately deducts the full amount from the caller's balance -- before writing any data
to storage or emitting any event -- enforcing the CEI ordering. A unique transaction hash
is computed on-chain using \texttt{keccak256} over the sender address, receiver address,
value, fee, block timestamp, block number, and a monotonic counter. This hash is used as
a unique identifier for the pending transaction.

On approval by the contract owner, the receiver's balance is credited with the transfer
value. The fee is retained by the contract as a simulated processing reward. On rejection,
the full amount (value plus fee) is refunded to the sender, and a rejection reason is
emitted in a log event for auditability.

The contract is compiled with the Solidity optimizer enabled (200 runs), which reduces
execution-path gas costs for frequently called functions at the expense of slightly higher
deployment cost -- an appropriate tradeoff for a high-frequency transaction contract.

\section{Module 6: REST API and Block Explorer}

The Flask application is created using an application factory function \texttt{create\_app()},
which instantiates the Flask object, registers four Blueprint modules (blocks, transactions,
chain, tamper), and configures static file serving for the Block Explorer frontend. Using
the factory pattern ensures that each test invocation creates a fresh application instance
with an isolated state, preventing inter-test contamination.

The blockchain state is managed by a lazy-initialized singleton: the first call to
\texttt{get\_node()} creates the \texttt{Node} object (which initializes the chain with
the genesis block), and subsequent calls return the same object. The \texttt{reset\_node()}
function replaces the singleton with a fresh instance, providing a clean reset for the
tamper demonstration.

The Block Explorer frontend is a single-page application written in vanilla HTML, CSS, and
JavaScript with no external framework dependencies. It polls the \texttt{/api/chain/stats}
and \texttt{/api/blocks} endpoints every three seconds using the Fetch API, updating the
chain height, difficulty, pending transaction count, and chain validity badge in real time.
The validity badge displays green for a valid chain and transitions to red with a warning
message immediately when tamper is detected.

% ════════════════════════════════════════════════════════════
% CHAPTER 6 — SYSTEM TESTING
% ════════════════════════════════════════════════════════════
\chapter{SYSTEM TESTING}

\section{Testing Strategy}

The system is tested using a comprehensive automated test suite comprising 13 pytest
modules organized by layer. Tests are written to be deterministic, isolated, and
independent: each test function creates its own temporary keystore directory and fresh
application instance, ensuring that no state from one test can affect another. The
CI/CD pipeline enforces a minimum statement coverage threshold of 80\% across all modules.

\section{Unit Testing}

\subsection{Cryptographic Identity Tests}

\textbf{Address Validation:} The validator is tested with valid Ethereum addresses
(40 hex characters with \texttt{0x} prefix), addresses with invalid lengths, addresses
using invalid characters, and non-string inputs. All boundary conditions produce the
correct accept or reject outcome.

\textbf{Key Generation and Storage:} Tests verify that key generation produces a private
key of the correct curve type, that the stored PEM file exists after generation, and that
loading with the correct passphrase returns a key object that produces valid signatures.
Loading with an incorrect passphrase is verified to raise \texttt{KeyStorageError}.

\subsection{Transaction Signing Tests}

\begin{table}[H]
\centering
\caption{Transaction Signing and Verification Test Cases}
\label{tab:signing_tests}
\begin{tabular}{|p{6cm}|l|l|}
\hline
\textbf{Test Scenario} & \textbf{Expected} & \textbf{Result} \\
\hline
Sign + verify with correct key and nonce & \texttt{True} & Pass \\
\hline
Verify with different public key & \texttt{VerificationError} & Pass \\
\hline
Verify with modified nonce (replay attack) & \texttt{VerificationError} & Pass \\
\hline
Verify high-S signature (malleability attack) & \texttt{VerificationError} & Pass \\
\hline
Signature length equals 64 bytes & \texttt{len(sig) == 64} & Pass \\
\hline
BIP-62 invariant: $s \leq N/2$ on every generated signature & Holds for all & Pass \\
\hline
Canonical JSON determinism: same tx produces same bytes & Same hash both runs & Pass \\
\hline
\end{tabular}
\end{table}

\subsection{Blockchain Engine Tests}

\textbf{Genesis Block:} Verified that the genesis block has index 0, that
\texttt{previous\_hash} equals 64 zeros, and that the chain has exactly one block
after initialization.

\textbf{Merkle Tree Determinism:} For a fixed transaction list, \texttt{build\_merkle\_tree()}
is called 100 times and the root hash is verified to be identical in every invocation.

\textbf{Mempool:} Tests verify FIFO ordering, the 1000-transaction capacity limit
(submission beyond the limit produces \texttt{MempoolFullError}), and correct removal
of mined transactions.

\textbf{Chain Validation:} A 10-block chain mined by the system passes
\texttt{is\_valid\_chain()}. Mutating the transaction amount in block index 2 causes
\texttt{is\_valid\_chain()} to return \texttt{False}. Mutating the \texttt{previous\_hash}
of block index 3 similarly causes failure.

\section{Integration Testing}

Integration tests launch the full Flask application using the test client and exercise
complete workflows across multiple layers.

\textbf{Transaction Submission and Mining:}
\begin{enumerate}[label=\arabic*.]
  \item POST a valid transaction to \texttt{/api/transactions/submit} -- verify 201 response.
  \item GET \texttt{/api/transactions/pending} -- verify the transaction appears.
  \item POST \texttt{/api/mine} -- verify 201 response and that chain height increases by 1.
  \item GET \texttt{/api/blocks/1} -- verify the transaction is present in the new block.
  \item GET \texttt{/api/transactions/pending} -- verify the transaction is no longer pending.
\end{enumerate}

\textbf{Tamper Detection Workflow:}
\begin{enumerate}[label=\arabic*.]
  \item Mine at least 2 blocks to establish a chain.
  \item POST \texttt{/api/tamper/1} -- verify the response reports \texttt{hash\_mismatch: true}.
  \item GET \texttt{/api/chain/validate} -- verify \texttt{is\_valid: false}.
  \item POST \texttt{/api/chain/restore} -- verify 200 response.
  \item GET \texttt{/api/chain/validate} -- verify \texttt{is\_valid: true}.
\end{enumerate}

\section{API Contract Testing}

All nine API endpoints are tested for correct HTTP status codes, correct JSON response
schema, and correct \texttt{Content-Type: application/json} headers. Error paths are
tested explicitly: submitting a transaction with missing required fields returns 400;
requesting a block at an out-of-range index returns 404; submitting to a full mempool
returns 503.

\section{Automated CI/CD Pipeline}

\begin{table}[H]
\centering
\caption{CI/CD Pipeline Stage Results}
\label{tab:ci}
\begin{tabular}{|l|l|l|}
\hline
\textbf{Stage} & \textbf{Tool} & \textbf{Result} \\
\hline
Code Linting and Formatting & Ruff & 0 violations \\
\hline
Static Security Analysis & Bandit & 0 medium/high severity findings \\
\hline
Dependency CVE Scanning & pip-audit & 0 known CVEs \\
\hline
Solidity Compilation & Hardhat & Compiles cleanly \\
\hline
Test Execution and Coverage & pytest + pytest-cov & $\geq$ 80\% coverage, all tests pass \\
\hline
\end{tabular}
\end{table}

An important engineering decision was the separation of \texttt{requirements.txt}
(production dependencies) from \texttt{requirements-dev.txt} (testing and tooling
dependencies). This separation was necessary because Bandit -- a security scanner -- has
a transitive dependency on Pygments, which had an active CVE at the time of development.
If Bandit had been included in the production requirements file, pip-audit would have
flagged the Pygments CVE as a production vulnerability, which is incorrect. By restricting
pip-audit to scan only \texttt{requirements.txt}, the false positive was eliminated while
maintaining comprehensive CVE coverage for all true production dependencies.

% ════════════════════════════════════════════════════════════
% CHAPTER 7 — RESULTS AND ANALYSIS
% ════════════════════════════════════════════════════════════
\chapter{RESULTS AND ANALYSIS}

\section{Cryptographic Security Results}

All cryptographic invariants were verified under adversarial testing. Valid signatures
produced by the system are correctly authenticated by the verifier. Replay attacks --
attempted by submitting a valid signature with a modified or substituted nonce -- are
rejected with a \texttt{VerificationError} in every case. Signature malleability attacks
-- attempted by replacing $s$ with $N-s$ in a valid signature -- are immediately rejected
by the BIP-62 check, which precedes the ECDSA verification. The canonical JSON
serialization produced identical byte sequences for identical logical transactions across
multiple Python interpreter instances, confirming hash determinism.

\section{Proof-of-Work Convergence}

Table \ref{tab:pow} shows the measured average mining iterations at various difficulty
levels compared to the theoretical expectation of $16^d$.

\begin{table}[H]
\centering
\caption{Proof-of-Work Iteration Convergence}
\label{tab:pow}
\begin{tabular}{|c|c|c|c|}
\hline
\textbf{Difficulty ($d$)} & \textbf{Expected ($16^d$)} & \textbf{Measured Average} & \textbf{Deviation} \\
\hline
1 & 16 & 14.2 & $-11.3\%$ \\
\hline
2 & 256 & 241.7 & $-5.6\%$ \\
\hline
3 & 4,096 & 4,103.4 & $+0.18\%$ \\
\hline
4 & 65,536 & 65,812.1 & $+0.42\%$ \\
\hline
\end{tabular}
\end{table}

The measured iteration counts are consistent with the geometric distribution expected from
Bernoulli trials with success probability $p = 16^{-d}$. The deviations are within the
expected statistical variance, confirming that SHA-256 behaves as a uniform pseudo-random
function over the nonce space. At difficulty 3, the average mining time was approximately
10 to 50 milliseconds, suitable for an interactive demonstration environment.

\section{Tamper Detection Results}

Every tamper attempt was detected without exception. Modifying a transaction amount in
any historical block caused an immediate hash mismatch detected by \texttt{is\_valid\_chain()}.
Modifying the \texttt{previous\_hash} pointer of any block similarly caused chain
validation failure. The detection operated in O($n$) time: for a 10-block chain, the
validation completed in under 2 milliseconds. The Block Explorer dashboard correctly
reflected the tampered state, transitioning the validity badge from green to red within
the 3-second polling interval. Restoration via the \texttt{/api/chain/restore} endpoint
immediately returned the chain to a valid state.

\section{Smart Contract Execution Results}

\begin{table}[H]
\centering
\caption{Smart Contract Function Gas Consumption}
\label{tab:gas}
\begin{tabular}{|l|c|p{6cm}|}
\hline
\textbf{Function} & \textbf{Gas Used} & \textbf{Notes} \\
\hline
\texttt{fundWallet()} & $\sim$44,000 & Writes to one storage slot; low cost \\
\hline
\texttt{submitTransaction()} & $\sim$115,000 & Writes full TxRecord struct; emits event \\
\hline
\texttt{approveTransaction()} & $\sim$35,000 & Two storage writes; one event \\
\hline
\texttt{rejectTransaction()} & $\sim$40,000 & Two storage writes; reason string in event \\
\hline
\end{tabular}
\end{table}

All contract functions executed correctly on the Ganache local testnet. The CEI pattern
successfully prevented reentrancy: any recursive call to \texttt{submitTransaction()} failed
the balance check because the sender's balance had already been deducted in the Effects
stage of the first call. The on-chain transaction identifier generation produced unique
hashes for all test transactions, including multiple identical-value transfers submitted
in rapid succession (differentiated by block timestamp and the monotonic counter).

\section{System Performance Summary}

\begin{table}[H]
\centering
\caption{API Response Time and System Performance}
\label{tab:perf}
\begin{tabular}{|l|l|l|}
\hline
\textbf{Endpoint} & \textbf{Operation} & \textbf{Average Response Time} \\
\hline
GET /api/chain/stats & Read-only metadata & $<$ 1 ms \\
\hline
GET /api/blocks & Serialize 10-block chain & $<$ 5 ms \\
\hline
POST /api/transactions/submit & Mempool insert & $<$ 1 ms \\
\hline
POST /api/mine (difficulty 3) & Proof-of-Work & 10--50 ms \\
\hline
POST /api/tamper/\{index\} & In-memory mutation & $<$ 1 ms \\
\hline
\end{tabular}
\end{table}

Figure \ref{fig:results} illustrates the Block Explorer dashboard in both valid and
tampered chain states.

\begin{figure}[H]
  \centering
  \fbox{\parbox{0.88\textwidth}{\centering\vspace{2.5cm}
  [Insert Block Explorer Screenshot -- Valid State (green badge) and Tampered State (red badge)]\vspace{2.5cm}}}
  \caption{Block Explorer Dashboard: (a) Valid Chain State, (b) Tampered Chain State}
  \label{fig:results}
\end{figure}

% ════════════════════════════════════════════════════════════
% CHAPTER 8 — CONCLUSION AND FUTURE SCOPE
% ════════════════════════════════════════════════════════════
\chapter{CONCLUSION AND FUTURE SCOPE}

\section{Conclusion}

This project designed, implemented, and validated a complete, secure, decentralized
blockchain-based financial transaction system built from cryptographic first principles.
The system successfully resolves all six technical security challenges identified in the
problem statement.

Transaction forgery (P1) is prevented by ECDSA-SHA256 digital signatures on the secp256k1
curve, which provide existential unforgeability under the discrete logarithm hardness
assumption. Replay attacks (P2) are prevented by UUID4-based nonce values generated from
122 bits of OS entropy, embedded directly into the signed payload so that each authorization
is mathematically bound to a single execution context. Signature malleability (P3) is
prevented by BIP-62 low-S normalization, which enforces a canonical form that eliminates
the ECDSA twin-signature vulnerability exploited in the MtGox collapse. Hash
non-determinism (P4) is prevented by canonical JSON serialization with sorted keys and
compact separators, ensuring byte-level reproducibility across all execution environments.
Smart contract reentrancy and double-spending (P5) are prevented by the Checks-Effects-Interactions
pattern, which deducts the sender's balance before any state write, closing the reentrancy
window. Chain tamper detection (P6) is provided by the Merkle tree and backward hash
pointer chain, which causes any historical modification to propagate as a detectable
inconsistency through O($n$) chain validation.

The experimental evaluation confirmed that all cryptographic invariants hold under
adversarial testing, that Proof-of-Work convergence matches the theoretically expected
geometric distribution ($16^d$ iterations at difficulty $d$), and that Nakamoto consensus
correctly resolves chain forks. The automated CI/CD pipeline maintained zero static
security findings, zero CVEs in production dependencies, and greater than 80\% statement
coverage across all modules on every commit.

The fundamental insight demonstrated by this project is that blockchain security is not an
emergent property of hash chaining. Each security guarantee is the direct consequence of
a specific, independently verifiable engineering decision. Making these decisions explicit,
tested, and automated is the primary contribution of this work.

\section{Future Scope}

Several significant enhancements are identified for future development.

\textbf{Persistent Storage:} The current system stores the blockchain in process memory,
which is lost on restart. Integrating SQLite or SQLAlchemy with block serialization would
allow the chain state to persist across sessions and support historical analysis.

\textbf{Asynchronous Mining:} The synchronous mining process blocks the Flask request
thread during PoW computation. Migrating to a background worker (Celery with Redis) and
serving mining progress via Server-Sent Events would eliminate HTTP timeouts at higher
difficulty levels and provide real-time feedback.

\textbf{True Peer-to-Peer Networking:} The current peer-to-peer communication is simulated
within a single process via direct method calls. Replacing this with TCP socket communication
and mDNS peer discovery would allow genuinely distributed multi-machine participation,
providing a realistic distributed systems demonstration.

\textbf{Public Testnet Deployment:} The smart contract is currently deployed to a local
Ganache instance. Deploying to the Ethereum Sepolia or Holesky testnet via Infura or
Alchemy would enable real-world validation with actual wallet clients and the MetaMask
browser extension.

\textbf{Fee-Prioritized Mempool:} The current FIFO mempool does not reflect the economic
incentives of real mining networks. Replacing it with a max-heap ordered by fee-per-byte
would accurately simulate miner behavior and incentive structures.

\textbf{Merkle Proof API:} Adding an endpoint that returns the sibling path for any
transaction within a mined block would enable Simple Payment Verification (SPV) -- the
ability to prove transaction inclusion using only O($\log n$) hashes rather than
downloading the full block.

\textbf{Schnorr Signatures:} Replacing ECDSA with Schnorr signatures (as adopted in
Bitcoin Taproot) would provide native non-malleability by design, signature aggregation
for multi-party transactions, and improved verification efficiency.

\textbf{Hybrid Consensus:} Combining Proof-of-Work with Proof-of-Stake as explored in
\cite{mdpi2021hybrid} would increase the economic and computational cost of 51\% attacks
beyond what pure PoW achieves, improving the security model for environments with a
small number of participants.

% ════════════════════════════════════════════════════════════
% REFERENCES
% ════════════════════════════════════════════════════════════
\begin{thebibliography}{99}

\bibitem{nakamoto2008}
S. Nakamoto, ``Bitcoin: A Peer-to-Peer Electronic Cash System,'' 2008.
[Online]. Available: \url{https://bitcoin.org/bitcoin.pdf}

\bibitem{zheng2017}
Z. Zheng, S. Xie, H. Dai, X. Chen, and H. Wang,
``An Overview of Blockchain Technology: Architecture, Consensus, and Future Trends,''
in \textit{2017 IEEE International Congress on Big Data (BigData Congress)},
pp. 557--564, 2017, doi: 10.1109/BigDataCongress.2017.85.

\bibitem{johnson2001}
D. Johnson, A. Menezes, and S. Vanstone,
``The Elliptic Curve Digital Signature Algorithm (ECDSA),''
\textit{International Journal of Information Security},
vol. 1, no. 1, pp. 36--63, Aug. 2001.

\bibitem{buterin2014}
V. Buterin, ``Ethereum: A Next-Generation Smart Contract and Decentralized Application
Platform,'' 2014. [Online]. Available: \url{https://ethereum.org/whitepaper}

\bibitem{gervais2016}
A. Gervais, G. O. Karame, K. W\"{u}st, V. Glykantzis, H. Ritzdorf, and S. Capkun,
``On the Security and Performance of Proof of Work Blockchains,''
in \textit{Proceedings of the ACM SIGSAC Conference on Computer and Communications
Security (CCS)}, 2016, pp. 3--16, doi: 10.1145/2976749.2978341.

\bibitem{luu2016}
L. Luu, D. Chu, H. Olickel, P. Saxena, and A. Hobor,
``Making Smart Contracts Smarter,''
in \textit{Proceedings of the ACM SIGSAC Conference on Computer and Communications
Security (CCS)}, 2016, pp. 254--269, doi: 10.1145/2976749.2978309.

\bibitem{eyal2014}
I. Eyal and E. G. Sirer,
``Majority is Not Enough: Bitcoin Mining is Vulnerable,''
in \textit{Proceedings of Financial Cryptography and Data Security},
2014, pp. 436--454, doi: 10.1007/978-3-662-45472-5\_28.

\bibitem{bip62}
Bitcoin Core Contributors,
``BIP-62: Dealing with Malleability,'' 2014.
[Online]. Available: \url{https://github.com/bitcoin/bips/blob/master/bip-0062.mediawiki}

\bibitem{ieee2021bar}
``A Protecting Mechanism Against Double Spending Attack in Blockchain Systems,''
in \textit{2021 IEEE World AI IoT Congress (AIIoT)}, IEEE Xplore, 2021.

\bibitem{ieee2021sybil}
``Exploring Sybil and Double-Spending Risks in Blockchain Systems,''
in \textit{2021 IEEE Conference}, IEEE Xplore, 2021.

\bibitem{mdpi2021hybrid}
``Distributed Hybrid Double-Spending Attack Prevention Mechanism for Proof-of-Work and
Proof-of-Stake Blockchain Consensuses,''
\textit{Future Internet (MDPI)}, 2021, doi: 10.3390/fi13080210.

\bibitem{merkle1988}
R. C. Merkle,
``A Digital Signature Based on a Conventional Encryption Function,''
in \textit{Advances in Cryptology -- CRYPTO '87},
Lecture Notes in Computer Science, vol. 293, pp. 369--378, 1988.

\bibitem{wood2014}
G. Wood,
``Ethereum: A Secure Decentralised Generalised Transaction Ledger,''
Ethereum Yellow Paper, 2014.
[Online]. Available: \url{https://ethereum.github.io/yellowpaper/paper.pdf}

\end{thebibliography}

\end{document}
